% Options for packages loaded elsewhere
% Options for packages loaded elsewhere
\PassOptionsToPackage{unicode}{hyperref}
\PassOptionsToPackage{hyphens}{url}
%
\documentclass[
  ignorenonframetext,
  aspectratio=169,
  russian,
]{beamer}
\newif\ifbibliography
\usepackage{pgfpages}
\setbeamertemplate{caption}[numbered]
\setbeamertemplate{caption label separator}{: }
\setbeamercolor{caption name}{fg=normal text.fg}
\beamertemplatenavigationsymbolshorizontal
% remove section numbering
\setbeamertemplate{part page}{
  \centering
  \begin{beamercolorbox}[sep=16pt,center]{part title}
    \usebeamerfont{part title}\insertpart\par
  \end{beamercolorbox}
}
\setbeamertemplate{section page}{
  \centering
  \begin{beamercolorbox}[sep=12pt,center]{section title}
    \usebeamerfont{section title}\insertsection\par
  \end{beamercolorbox}
}
\setbeamertemplate{subsection page}{
  \centering
  \begin{beamercolorbox}[sep=8pt,center]{subsection title}
    \usebeamerfont{subsection title}\insertsubsection\par
  \end{beamercolorbox}
}
% Prevent slide breaks in the middle of a paragraph
\widowpenalties 1 10000
\raggedbottom
\AtBeginPart{
  \frame{\partpage}
}
\AtBeginSection{
  \ifbibliography
  \else
    \frame{\sectionpage}
  \fi
}
\AtBeginSubsection{
  \frame{\subsectionpage}
}
\usepackage{iftex}
\ifPDFTeX
  \usepackage[T1]{fontenc}
  \usepackage[utf8]{inputenc}
  \usepackage{textcomp} % provide euro and other symbols
\else % if luatex or xetex
  \usepackage{unicode-math} % this also loads fontspec
  \defaultfontfeatures{Scale=MatchLowercase}
  \defaultfontfeatures[\rmfamily]{Ligatures=TeX,Scale=1}
\fi
\usepackage{lmodern}

\ifPDFTeX\else
  % xetex/luatex font selection
\fi
% Use upquote if available, for straight quotes in verbatim environments
\IfFileExists{upquote.sty}{\usepackage{upquote}}{}
\IfFileExists{microtype.sty}{% use microtype if available
  \usepackage[]{microtype}
  \UseMicrotypeSet[protrusion]{basicmath} % disable protrusion for tt fonts
}{}


\usepackage{longtable,booktabs,array}
\newcounter{none} % for unnumbered tables
\usepackage{calc} % for calculating minipage widths
\usepackage{caption}
% Make caption package work with longtable
\makeatletter
\def\fnum@table{\tablename~\thetable}
\makeatother
\usepackage{graphicx}
\makeatletter
\newsavebox\pandoc@box
\newcommand*\pandocbounded[1]{% scales image to fit in text height/width
  \sbox\pandoc@box{#1}%
  \Gscale@div\@tempa{\textheight}{\dimexpr\ht\pandoc@box+\dp\pandoc@box\relax}%
  \Gscale@div\@tempb{\linewidth}{\wd\pandoc@box}%
  \ifdim\@tempb\p@<\@tempa\p@\let\@tempa\@tempb\fi% select the smaller of both
  \ifdim\@tempa\p@<\p@\scalebox{\@tempa}{\usebox\pandoc@box}%
  \else\usebox{\pandoc@box}%
  \fi%
}
% Set default figure placement to htbp
\def\fps@figure{htbp}
\makeatother



\ifLuaTeX
\usepackage[bidi=basic,shorthands=off,provide=*]{babel}
\else
\usepackage[bidi=default,shorthands=off,provide=*]{babel}
\fi
\ifLuaTeX
  \usepackage{selnolig} % disable illegal ligatures
\fi


\setlength{\emergencystretch}{3em} % prevent overfull lines

\providecommand{\tightlist}{%
  \setlength{\itemsep}{0pt}\setlength{\parskip}{0pt}}



 

\usepackage[]{csquotes}

\IfFileExists{plex-otf.sty}{
  %% Full TeXlive
  % \usepackage[%
  %   % math,
  %   RM={Scale=0.94},SS={Scale=0.94},SScon={Scale=0.94},TT={Scale=MatchLowercase,FakeStretch=0.9},DefaultFeatures={Ligatures=Common}
  % ]{plex-otf}
}{
  %% TinyTeX
  \usepackage{libertine}
}

%%% Load theme
% https://deic.uab.cat/~iblanes/beamer_gallery/
\IfFileExists{beamerthemegotham.sty}{
  %% Full TeXlive
  \usetheme{gotham}
  \gothamset{
    numbering=totalpagenumber,
    parttocframe default=off,
    sectiontocframe default=off,
    subsectiontocframe default=off,
  }
}{
  %% TinyTeX
  \usetheme{Madrid}
}
\makeatletter
\@ifpackageloaded{caption}{}{\usepackage{caption}}
\AtBeginDocument{%
\ifdefined\contentsname
  \renewcommand*\contentsname{Содержание}
\else
  \newcommand\contentsname{Содержание}
\fi
\ifdefined\listfigurename
  \renewcommand*\listfigurename{Список иллюстраций}
\else
  \newcommand\listfigurename{Список иллюстраций}
\fi
\ifdefined\listtablename
  \renewcommand*\listtablename{Список таблиц}
\else
  \newcommand\listtablename{Список таблиц}
\fi
\ifdefined\figurename
  \renewcommand*\figurename{Рисунок}
\else
  \newcommand\figurename{Рисунок}
\fi
\ifdefined\tablename
  \renewcommand*\tablename{Таблица}
\else
  \newcommand\tablename{Таблица}
\fi
}
\@ifpackageloaded{float}{}{\usepackage{float}}
\floatstyle{ruled}
\@ifundefined{c@chapter}{\newfloat{codelisting}{h}{lop}}{\newfloat{codelisting}{h}{lop}[chapter]}
\floatname{codelisting}{Список}
\newcommand*\listoflistings{\listof{codelisting}{Листинги}}
\makeatother
\makeatletter
\makeatother
\makeatletter
\@ifpackageloaded{caption}{}{\usepackage{caption}}
\@ifpackageloaded{subcaption}{}{\usepackage{subcaption}}
\makeatother

\usepackage{bookmark}
\IfFileExists{xurl.sty}{\usepackage{xurl}}{} % add URL line breaks if available
\urlstyle{same}
\hypersetup{
  pdftitle={Лабораторная работа №4},
  pdfauthor={Лопатченко Полина Андреевна},
  pdflang={ru-RU},
  hidelinks,
  pdfcreator={LaTeX via pandoc}}


\title{\texorpdfstring{Лабораторная работа №4}{Лабораторная работа №4}}
\subtitle{\texorpdfstring{Продвинутое использование
git.}{Продвинутое использование git.}}
\author{Лопатченко Полина Андреевна}
\date{2026-03-07}

\begin{document}
\frame{\titlepage}

\renewcommand*\contentsname{Содержание}
\begin{frame}[allowframebreaks]
  \frametitle{Содержание}
  \setcounter{tocdepth}{1}
  \tableofcontents
\end{frame}
\setcounter{tocdepth}{1}
\tableofcontents
}

\section{Информация}\label{ux438ux43dux444ux43eux440ux43cux430ux446ux438ux44f}

\begin{frame}{Докладчик}
\protect\phantomsection\label{ux434ux43eux43aux43bux430ux434ux447ux438ux43a}
\begin{columns}[c]
\begin{column}{0.7\linewidth}
\begin{itemize}[<+->]
\tightlist
\item
  Лпатченко Полина Андреевна
\item
  студент
\item
  НКАбд-04-25
\item
  Российский университет дружбы народов им. П. Лумумбы
\item
  \href{mailro:1032253529@rudn.ru}{1032253529@rudn.ru}
\item
  \url{https://PALopatchenko.github.io/ru/}
\end{itemize}
\end{column}

\begin{column}{0.3\linewidth}
\pandocbounded{
\includegraphics[keepaspectratio]{./image/я.jpg}}
\end{column}
\end{columns}
\end{frame}

\section{Вводная
часть}\label{ux432ux432ux43eux434ux43dux430ux44f-ux447ux430ux441ux442ux44c}

\begin{frame}{Актуальность}
\protect\phantomsection\label{ux430ux43aux442ux443ux430ux43bux44cux43dux43eux441ux442ux44c}
\begin{itemize}[<+->]
\tightlist
\item
  Важно донести результаты своих исследований до окружающих
\item
  Научная презентация --- рабочий инструмент исследователя
\item
  Необходимо создавать презентацию быстро
\item
  Желательна минимизация усилий для создания презентации
\end{itemize}
\end{frame}

\begin{frame}{Объект и предмет исследования}
\protect\phantomsection\label{ux43eux431ux44aux435ux43aux442-ux438-ux43fux440ux435ux434ux43cux435ux442-ux438ux441ux441ux43bux435ux434ux43eux432ux430ux43dux438ux44f}
\begin{itemize}[<+->]
\tightlist
\item
  Презентация как текст
\item
  Программное обеспечение для создания презентаций
\item
  Входные и выходные форматы презентаций
\end{itemize}
\end{frame}

\begin{frame}{Цели и задачи}
\protect\phantomsection\label{ux446ux435ux43bux438-ux438-ux437ux430ux434ux430ux447ux438}
\end{frame}

\section{Цели}\label{ux446ux435ux43bux438}

\section{Задание}\label{ux437ux430ux434ux430ux43dux438ux435}

\begin{frame}{Задание}
\begin{itemize}[<+->]
\tightlist
\item
  Выполнить работу для тестового репозитория.
\item
  Преобразовать рабочий репозиторий в репозиторий с git-flow и
  conventional commits.
\end{itemize}
\end{frame}

\section{Теоретическое
введение}\label{ux442ux435ux43eux440ux435ux442ux438ux447ux435ux441ux43aux43eux435-ux432ux432ux435ux434ux435ux43dux438ux435}

\begin{frame}[fragile]{Рабочий процесс Gitflow}
\protect\phantomsection\label{ux440ux430ux431ux43eux447ux438ux439-ux43fux440ux43eux446ux435ux441ux441-gitflow}
\begin{itemize}[<+->]
\tightlist
\item
  Рабочий процесс \emph{Gitflow Workflow}. Будем описывать его с
  использованием пакета \texttt{git-flow}.
\end{itemize}

\begin{block}{Общая информация}
\protect\phantomsection\label{ux43eux431ux449ux430ux44f-ux438ux43dux444ux43eux440ux43cux430ux446ux438ux44f}
\begin{itemize}[<+->]
\tightlist
\item
  Gitflow Workflow опубликована и популяризована Винсентом Дриссеном.
\item
  Gitflow Workflow предполагает выстраивание строгой модели ветвления с
  учётом выпуска проекта.
\item
  Данная модель отлично подходит для организации рабочего процесса на
  основе релизов.
\item
  Работа по модели Gitflow включает создание отдельной ветки для
  исправлений ошибок в рабочей среде.
\item
  Последовательность действий при работе по модели Gitflow:

  \begin{itemize}[<+->]
  \tightlist
  \item
    Из ветки \texttt{master} создаётся ветка \texttt{develop}.
  \item
    Из ветки \texttt{develop} создаётся ветка \texttt{release}.
  \item
    Из ветки \texttt{develop} создаются ветки \texttt{feature}.
  \item
    Когда работа над веткой \texttt{feature} завершена, она сливается с
    веткой \texttt{develop}.
  \item
    Когда работа над веткой релиза \texttt{release} завершена, она
    сливается в ветки \texttt{develop} и \texttt{master}.
  \item
    Если в \texttt{master} обнаружена проблема, из \texttt{master}
    создаётся ветка \texttt{hotfix}.
  \item
    Когда работа над веткой исправления \texttt{hotfix} завершена, она
    сливается в ветки \texttt{develop} и \texttt{master}.
  \end{itemize}
\end{itemize}
\end{block}

\begin{block}{Процесс работы с Gitflow}
\protect\phantomsection\label{ux43fux440ux43eux446ux435ux441ux441-ux440ux430ux431ux43eux442ux44b-ux441-gitflow}
\begin{enumerate}[<+->]
\item
  Основные ветки (\texttt{master}) и ветки разработки (\texttt{develop})

  \begin{itemize}[<+->]
  \tightlist
  \item
    Для фиксации истории проекта в рамках этого процесса вместо одной
    ветки \texttt{master} используются две ветки. В ветке
    \texttt{master} хранится официальная история релиза, а ветка
    \texttt{develop} предназначена для объединения всех функций. Кроме
    того, для удобства рекомендуется присваивать всем коммитам в ветке
    \texttt{master} номер версии.
  \item
    При использовании библиотеки расширений \texttt{git-flow} нужно
    инициализировать структуру в существующем репозитории:
  \end{itemize}

\begin{verbatim}
git flow init
\end{verbatim}

  \begin{itemize}[<+->]
  \tightlist
  \item
    Для github параметр \texttt{Version\ tag\ prefix} следует установить
    в \texttt{v}.
  \item
    После этого проверьте, на какой ветке Вы находитесь:
  \end{itemize}

\begin{verbatim}
git branch
\end{verbatim}
\item
  Функциональные ветки (feature)

  \begin{itemize}[<+->]
  \tightlist
  \item
    Под каждую новую функцию должна быть отведена собственная ветка,
    которую можно отправлять в центральный репозиторий для создания
    резервной копии или совместной работы команды. Ветки
    \texttt{feature} создаются не на основе \texttt{master}, а на основе
    \texttt{develop}. Когда работа над функцией завершается,
    соответствующая ветка сливается обратно с веткой \texttt{develop}.
    Функции не следует отправлять напрямую в ветку \texttt{master}.
  \item
    Как правило, ветки \texttt{feature} создаются на основе последней
    ветки \texttt{develop}.

    \begin{enumerate}[<+->]
    \item
      Создание функциональной ветки

      \begin{itemize}[<+->]
      \tightlist
      \item
        Создадим новую функциональную ветку:
      \end{itemize}

\begin{verbatim}
git flow feature start feature_branch
\end{verbatim}

      \begin{itemize}[<+->]
      \tightlist
      \item
        Создадим новую функциональную ветку:
      \end{itemize}

\begin{verbatim}
git flow feature start feature_branch
\end{verbatim}

      \begin{itemize}[<+->]
      \tightlist
      \item
        Далее работаем как обычно.
      \end{itemize}
    \item
      Окончание работы с функциональной веткой

      \begin{itemize}[<+->]
      \tightlist
      \item
        По завершении работы над функцией следует объединить ветку
        \texttt{feature\_branch} с \texttt{develop}:
      \end{itemize}

\begin{verbatim}
git flow feature finish feature_branch
\end{verbatim}
    \item
      Ветки выпуска (release)

      \begin{itemize}[<+->]
      \tightlist
      \item
        Когда в ветке \texttt{develop} оказывается достаточно функций
        для выпуска, из ветки \texttt{develop} создаётся ветка
        \texttt{release}. Создание этой ветки запускает следующий цикл
        выпуска, и с этого момента новые функции добавить больше нельзя
        --- допускается лишь отладка, создание документации и решение
        других задач. Когда подготовка релиза завершается, ветка
        \texttt{release} сливается с \texttt{master} и ей присваивается
        номер версии. После нужно выполнить слияние с веткой
        \texttt{develop}, в которой с момента создания ветки релиза
        могли возникнуть изменения.
      \item
        Благодаря тому, что для подготовки выпусков используется
        специальная ветка, одна команда может дорабатывать текущий
        выпуск, в то время как другая команда продолжает работу над
        функциями для следующего.
      \item
        Создать новую ветку release можно с помощью следующей команды:
      \end{itemize}

\begin{verbatim}
git flow release start 1.0.0
\end{verbatim}

      \begin{itemize}[<+->]
      \tightlist
      \item
        Для завершения работы на ветке release используются следующие
        команды:
      \end{itemize}

\begin{verbatim}
git flow release finish 1.0.0
\end{verbatim}
    \item
      Ветки исправления (hotfix)

      \begin{itemize}[<+->]
      \tightlist
      \item
        Ветки поддержки или ветки \texttt{hotfix} используются для
        быстрого внесения исправлений в рабочие релизы. Они создаются от
        ветки \texttt{master}. Это единственная ветка, которая должна
        быть создана непосредственно от \texttt{master}. Как только
        исправление завершено, ветку следует объединить с
        \texttt{master} и \texttt{develop}. Ветка \texttt{master} должна
        быть помечена обновлённым номером версии.
      \item
        Наличие специальной ветки для исправления ошибок позволяет
        команде решать проблемы, не прерывая остальную часть рабочего
        процесса и не ожидая следующего цикла релиза.
      \item
        Ветку \texttt{hotfix} можно создать с помощью следующих команд:
      \end{itemize}

\begin{verbatim}
git flow hotfix start hotfix_branch
\end{verbatim}

      По завершении работы ветка \texttt{hotfix} объединяется с
      \texttt{master} и \texttt{develop}:
      \texttt{git\ flow\ hotfix\ finish\ hotfix\_branch}
    \end{enumerate}
  \end{itemize}
\end{enumerate}
\end{block}
\end{frame}

\begin{frame}[fragile]{Семантическое версионирование}
\protect\phantomsection\label{ux441ux435ux43cux430ux43dux442ux438ux447ux435ux441ux43aux43eux435-ux432ux435ux440ux441ux438ux43eux43dux438ux440ux43eux432ux430ux43dux438ux435}
Семантический подход в версионированию программного обеспечения.

\begin{block}{Краткое описание семантического версионирования}
\protect\phantomsection\label{ux43aux440ux430ux442ux43aux43eux435-ux43eux43fux438ux441ux430ux43dux438ux435-ux441ux435ux43cux430ux43dux442ux438ux447ux435ux441ux43aux43eux433ux43e-ux432ux435ux440ux441ux438ux43eux43dux438ux440ux43eux432ux430ux43dux438ux44f}
\begin{itemize}[<+->]
\item
  Семантическое версионирование описывается в манифесте семантического
  версионирования.
\item
  Кратко его можно описать следующим образом:

  \begin{itemize}[<+->]
  \tightlist
  \item
    Версия задаётся в виде кортежа
    \texttt{МАЖОРНАЯ\_ВЕРСИЯ.МИНОРНАЯ\_ВЕРСИЯ.ПАТЧ}.
  \item
    Номер версии следует увеличивать:

    \begin{itemize}[<+->]
    \tightlist
    \item
      МАЖОРНУЮ версию, когда сделаны обратно несовместимые изменения
      API.
    \item
      МИНОРНУЮ версию, когда вы добавляете новую функциональность, не
      нарушая обратной совместимости.
    \item
      ПАТЧ-версию, когда вы делаете обратно совместимые исправления.
    \end{itemize}
  \item
    Дополнительные обозначения для предрелизных и билд-метаданных
    возможны как дополнения к МАЖОРНАЯ.МИНОРНАЯ.ПАТЧ формату.
  \end{itemize}
\end{itemize}
\end{block}

\begin{block}{Программное обеспечение}
\protect\phantomsection\label{ux43fux440ux43eux433ux440ux430ux43cux43cux43dux43eux435-ux43eux431ux435ux441ux43fux435ux447ux435ux43dux438ux435}
\begin{itemize}[<+->]
\item
  Для реализации семантического версионирования создано несколько
  программных продуктов.
\item
  При этом лучше всего использовать комплексные продукты, которые
  используют информацию из коммитов системы версионирования.
\item
  Коммиты должны иметь стандартизованный вид.
\item
  В семантическое версионирование применяется вместе с
  \emph{общепринятыми коммитами}.

  \begin{enumerate}[<+->]
  \item
    Пакет Conventional Changelog

    \begin{itemize}[<+->]
    \tightlist
    \item
      Пакет Conventional Changelog является комплексным решением по
      управлению коммитами и генерации журнала изменений.
    \item
      Содержит набор утилит, которые можно использовать по-отдельности.
    \end{itemize}
  \end{enumerate}
\end{itemize}
\end{block}
\end{frame}

\begin{frame}[fragile]{Общепринятые коммиты}
\protect\phantomsection\label{ux43eux431ux449ux435ux43fux440ux438ux43dux44fux442ux44bux435-ux43aux43eux43cux43cux438ux442ux44b}
Использование спецификации Conventional Commits.

\begin{block}{Описание}
\protect\phantomsection\label{ux43eux43fux438ux441ux430ux43dux438ux435}
Спецификация Conventional Commits:

\begin{itemize}[<+->]
\item
  Соглашение о том, как нужно писать сообщения commit'ов.
\item
  Совместимо с SemVer. Даже вернее сказать, сильно связано с
  семантическим версионированием.
\item
  Регламентирует структуру и основные типы коммитов.

  \begin{enumerate}[<+->]
  \item
    Структура коммита ``` ():

    Или, по-русски:

    (): {[}необязательное тело{]} {[}необязательный нижний колонтитул{]}
    ```

    \begin{itemize}[<+->]
    \tightlist
    \item
      Заголовок является обязательным.
    \item
      Любая строка сообщения о фиксации не может быть длиннее 100
      символов.
    \item
      Тема (subject) содержит краткое описание изменения.

      \begin{itemize}[<+->]
      \tightlist
      \item
        Используйте повелительное наклонение в настоящем времени:
        «изменить» (\enquote{change} not \enquote{changed} nor
        \enquote{changes}).
      \item
        Не используйте заглавную первую букву.
      \item
        Не ставьте точку в конце.
      \end{itemize}
    \item
      Тело (body) должно включать мотивацию к изменению и
      противопоставлять это предыдущему поведению.

      \begin{itemize}[<+->]
      \tightlist
      \item
        Как и в теме, используйте повелительное наклонение в настоящем
        времени.
      \end{itemize}
    \item
      Нижний колонтитул (footer) должен содержать любую информацию о
      критических изменениях.

      \begin{itemize}[<+->]
      \tightlist
      \item
        Следует использовать для указания внешних ссылок, контекста
        коммита или другой мета информации.
      \item
        Также содержит ссылку на issue (например, на github), который
        закрывает эта фиксация.
      \item
        Критические изменения должны начинаться со слова
        \texttt{BREAKING\ CHANGE}: с пробела или двух символов новой
        строки. Затем для этого используется остальная часть сообщения
        фиксации.
      \end{itemize}
    \end{itemize}
  \item
    Типы коммитов

    \begin{enumerate}[<+->]
    \tightlist
    \item
      Базовые типы коммитов

      \begin{itemize}[<+->]
      \tightlist
      \item
        \texttt{fix:} --- коммит типа fix исправляет ошибку (bug) в
        вашем коде (он соответствует PATCH в SemVer).
      \item
        \texttt{feat:} --- коммит типа feat добавляет новую функцию
        (feature) в ваш код (он соответствует MINOR в SemVer).
      \item
        \texttt{BREAKING\ CHANGE:} --- коммит, который содержит текст
        \texttt{BREAKING\ CHANGE:} в начале своего не обязательного тела
        сообщения (body) или в подвале (footer), добавляет изменения,
        нарушающие обратную совместимость вашего API (он соответствует
        MAJOR в SemVer). BREAKING CHANGE может быть частью коммита
        любого типа.
      \item
        \texttt{revert:} --- если фиксация отменяет предыдущую фиксацию.
        Начинается с revert:, за которым следует заголовок отменённой
        фиксации. В теле должно быть написано: Это отменяет фиксацию
        (это SHA-хэш отменяемой фиксации).
      \item
        Другое: коммиты с типами, которые отличаются от \texttt{fix:} и
        \texttt{feat:}, также разрешены. Например,
        @commitlint/config-conventional (основанный на The Angular
        convention) рекомендует: chore:, docs:, style:, refactor:,
        perf:, test:, и другие.
      \end{itemize}
    \item
      Соглашения The Angular convention

      \begin{itemize}[<+->]
      \tightlist
      \item
        Одно из популярных соглашений о поддержке исходных кодов ---
        конвенция Angular (The Angular convention).
      \end{itemize}

      \begin{enumerate}[<+->]
      \tightlist
      \item
        Типы коммитов The Angular convention Конвенция Angular (The
        Angular convention) требует следующие типы коммитов:
      \end{enumerate}

      \begin{itemize}[<+->]
      \tightlist
      \item
        \texttt{build:} --- изменения, влияющие на систему сборки или
        внешние зависимости (примеры областей (scope): gulp, broccoli,
        npm).
      \item
        \texttt{ci:} --- изменения в файлах конфигурации и скриптах CI
        (примеры областей: Travis, Circle, BrowserStack, SauceLabs).
      \item
        \texttt{docs:} --- изменения только в документации.
      \item
        \texttt{feat:} --- новая функция.
      \item
        \texttt{fix:} --- исправление ошибок.
      \item
        \texttt{perf:} --- изменение кода, улучшающее
        производительность.
      \item
        \texttt{refactor:} --- Изменение кода, которое не исправляет
        ошибку и не добавляет функции (рефакторинг кода).
      \item
        \texttt{style:} --- изменения, не влияющие на смысл кода
        (пробелы, форматирование, отсутствие точек с запятой и т. д.).
      \item
        \texttt{test:} --- добавление недостающих тестов или исправление
        существующих тестов.
      \end{itemize}

      \begin{enumerate}[<+->]
      \tightlist
      \item
        Области действия (scope) Областью действия должно быть имя
        затронутого пакета npm (как его воспринимает человек, читающий
        журнал изменений, созданный из сообщений фиксации). Есть
        несколько исключений из правила «использовать имя пакета»:
      \end{enumerate}

      \begin{itemize}[<+->]
      \tightlist
      \item
        \texttt{packaging} --- используется для изменений, которые
        изменяют структуру пакета, например, изменения общедоступного
        пути.
      \item
        \texttt{changelog} --- используется для обновления примечаний к
        выпуску в CHANGELOG.md.
      \item
        отсутствует область действия --- полезно для изменений стиля,
        тестирования и рефакторинга, которые выполняются во всех пакетах
        (например, style: добавить отсутствующие точки с запятой).
      \end{itemize}
    \item
      Соглашения @commitlint/config-conventional Соглашение
      @commitlint/config-conventional входит в пакет Conventional
      Changelog. В целом в этом соглашении придерживаются соглашения
      Angular.
    \end{enumerate}
  \end{enumerate}
\end{itemize}
\end{block}
\end{frame}

\section{Выполнение лабораторной
работы}\label{ux432ux44bux43fux43eux43bux43dux435ux43dux438ux435-ux43bux430ux431ux43eux440ux430ux442ux43eux440ux43dux43eux439-ux440ux430ux431ux43eux442ux44b}

\begin{frame}[fragile]{Установка программного обеспечения}
\protect\phantomsection\label{ux443ux441ux442ux430ux43dux43eux432ux43aux430-ux43fux440ux43eux433ux440ux430ux43cux43cux43dux43eux433ux43e-ux43eux431ux435ux441ux43fux435ux447ux435ux43dux438ux44f}
\begin{block}{Установка git-flow}
\protect\phantomsection\label{ux443ux441ux442ux430ux43dux43eux432ux43aux430-git-flow}
Установила git-flow (см. \hyperlink{fig-001}{рис.1}).

\begin{figure}

\centering{

\includegraphics[width=0.7\linewidth,height=\textheight,keepaspectratio]{image/1.png}

}

\caption{\label{fig-001}Установка git-flow}

\end{figure}%
\end{block}

\begin{block}{Установка Node.js}
\protect\phantomsection\label{ux443ux441ux442ux430ux43dux43eux432ux43aux430-node.js}
Установила Node.js (см. \hyperlink{fig-002}{рис.2}).

\begin{figure}

\centering{

\includegraphics[width=0.7\linewidth,height=\textheight,keepaspectratio]{image/2.png}

}

\caption{\label{fig-002}Установка Node.js}

\end{figure}%
\end{block}

\begin{block}{Настройка Node.js}
\protect\phantomsection\label{ux43dux430ux441ux442ux440ux43eux439ux43aux430-node.js}
Запустила \texttt{pnpm\ setup} и выполнил
\texttt{source\ \textasciitilde{}/.bashrc} (см.
\hyperlink{fig-003}{рис.3}).

\begin{figure}

\centering{

\includegraphics[width=0.7\linewidth,height=\textheight,keepaspectratio]{image/3.png}

}

\caption{\label{fig-003}Запуск pnpm setup и выполнение source
\textasciitilde/.bashrc}

\end{figure}%
\end{block}
\end{frame}

\begin{frame}[fragile]{Практический сценарий использования git}
\protect\phantomsection\label{ux43fux440ux430ux43aux442ux438ux447ux435ux441ux43aux438ux439-ux441ux446ux435ux43dux430ux440ux438ux439-ux438ux441ux43fux43eux43bux44cux437ux43eux432ux430ux43dux438ux44f-git}
\begin{block}{Создание репозитория git}
\protect\phantomsection\label{ux441ux43eux437ux434ux430ux43dux438ux435-ux440ux435ux43fux43eux437ux438ux442ux43eux440ux438ux44f-git}
\begin{enumerate}[<+->]
\tightlist
\item
  Подключение репозитория к github
\end{enumerate}

\begin{itemize}[<+->]
\tightlist
\item
  Создала репозиторий на GitHub. Назвал его \texttt{git-extended} (см.
  \hyperlink{fig-004}{рис.4}).
\end{itemize}

\begin{figure}

\centering{

\includegraphics[width=0.7\linewidth,height=\textheight,keepaspectratio]{image/4.png}

}

\caption{\label{fig-004}Создание репозитория GitHub}

\end{figure}%

\begin{itemize}[<+->]
\tightlist
\item
  Сделала первый коммит и выложил его на github (см.
  \hyperlink{fig-005}{рис.5}).
\end{itemize}

\begin{figure}

\centering{

\includegraphics[width=0.7\linewidth,height=\textheight,keepaspectratio]{image/5.png}

}

\caption{\label{fig-005}Создание первого коммита и его загрузка на
github}

\end{figure}%

\begin{enumerate}[<+->]
\tightlist
\item
  Конфигурация общепринятых коммитов (см. \hyperlink{fig-006}{рис.6}).
\end{enumerate}

\begin{itemize}[<+->]
\tightlist
\item
  Конфигурация пакетов Node.js (см. \hyperlink{fig-006}{рис.6}).
\end{itemize}

\begin{figure}

\centering{

\includegraphics[width=0.7\linewidth,height=\textheight,keepaspectratio]{image/6.png}

}

\caption{\label{fig-006}Выполнение конфигурации пакетов Node.js}

\end{figure}%

Заполнила несколько параметров пакета и сконфигурировал формат коммитов
(см. \hyperlink{fig-007}{рис.7}).

\begin{figure}

\centering{

\includegraphics[width=0.7\linewidth,height=\textheight,keepaspectratio]{image/7.png}

}

\caption{\label{fig-007}Заполнение нескольких параметров пакета и
конфигурация формата коммитов}

\end{figure}%

\begin{enumerate}[<+->]
\tightlist
\item
  Конфигурация git-flow
\end{enumerate}

\begin{itemize}[<+->]
\tightlist
\item
  Инициализировала git-flow (см. \hyperlink{fig-011}{рис.11}).
\end{itemize}

\begin{figure}

\centering{

\includegraphics[width=0.7\linewidth,height=\textheight,keepaspectratio]{image/11.png}

}

\caption{\label{fig-011}Инициализация git-flow}

\end{figure}%

\begin{itemize}[<+->]
\tightlist
\item
  Проверила, что я на ветке \texttt{develop}, загрузил весь репозиторий
  в хранилище, установил внешнюю ветку как вышестоящую для этой ветки и
  создал релиз с версией 1.0.0 (см. \hyperlink{fig-012}{рис.12}).
\end{itemize}

\begin{figure}

\centering{

\includegraphics[width=0.7\linewidth,height=\textheight,keepaspectratio]{image/12.png}

}

\caption{\label{fig-012}Проверка нахождения на ветке develop, загрузка
всего репозитория в хранилище, установка внешней ветки как вышестоящей
для этой ветки и создание релиза с версией 1.0.0}

\end{figure}%

Создала журнал изменений и добавила его в индекс (см.
\hyperlink{fig-013}{рис.13}).

\begin{figure}

\centering{

\includegraphics[width=0.7\linewidth,height=\textheight,keepaspectratio]{image/13.png}

}

\caption{\label{fig-013}Создание журнала изменений и добавление его в
индекс}

\end{figure}%
\end{block}

\begin{block}{Работа с репозиторием git}
\protect\phantomsection\label{ux440ux430ux431ux43eux442ux430-ux441-ux440ux435ux43fux43eux437ux438ux442ux43eux440ux438ux435ux43c-git}
\begin{enumerate}[<+->]
\item
  Разработка новой функциональности
\item
  Создание релиза git-flow
\end{enumerate}

\begin{itemize}[<+->]
\tightlist
\item
  Создала релиз с версией \texttt{1.2.3}, обновила номер версии в
  \texttt{package.json}, создал журнал изменений и добавил журнал
  изменений в индекс (см. \hyperlink{fig-016}{рис.16}).
\end{itemize}

\begin{figure}

\centering{

\includegraphics[width=0.7\linewidth,height=\textheight,keepaspectratio]{image/16.png}

}

\caption{\label{fig-016}Создание релиза с версией \texttt{1.2.3},
обновление номера версии в \texttt{package.json}, создание журнала
изменений и добавление журнал изменений в индекс}

\end{figure}%

\begin{itemize}[<+->]
\tightlist
\item
  Залила релизную ветку в основную ветку, отправила данные на github,
  создала релиз на github с комментариями из журнала изменений (см.
  \hyperlink{fig-017}{рис.17}).
\end{itemize}

\begin{figure}

\centering{

\includegraphics[width=0.7\linewidth,height=\textheight,keepaspectratio]{image/17.png}

}

\caption{\label{fig-017}Залил релизную ветку в основную ветку, отправка
данных на github, создание релиза на github с комментариями из журнала
изменений}

\end{figure}%
\end{block}
\end{frame}

\begin{frame}[fragile]{Преобразование рабочего репозитория в репозиторий
с git-flow и conventional commits}
\protect\phantomsection\label{ux43fux440ux435ux43eux431ux440ux430ux437ux43eux432ux430ux43dux438ux435-ux440ux430ux431ux43eux447ux435ux433ux43e-ux440ux435ux43fux43eux437ux438ux442ux43eux440ux438ux44f-ux432-ux440ux435ux43fux43eux437ux438ux442ux43eux440ux438ux439-ux441-git-flow-ux438-conventional-commits}
Зашла в папку \texttt{os-intro}, создала новый репозиторий, добавила все
файлы в индекс, сохранила первый коммит и создала \texttt{package.json}
(см. \hyperlink{fig-018}{рис.18}).

\begin{figure}

\centering{

\includegraphics[width=0.7\linewidth,height=\textheight,keepaspectratio]{image/18.png}

}

\caption{\label{fig-018}Переход в папку os-intro, создание нового
репозитория, добавление всех файлов в индекс, сохранение первого коммита
и создание package.json}

\end{figure}%

Заполнила несколько параметров пакета и сконфигурировал формат коммитов
(см. \hyperlink{fig-019}{рис.19}).

\begin{figure}

\centering{

\includegraphics[width=0.7\linewidth,height=\textheight,keepaspectratio]{image/19.png}

}

\caption{\label{fig-019}Заполнение параметров пакета и конфигурация
формата коммитов}

\end{figure}%

Сделала \texttt{git\ flow\ init}, настроил \texttt{master} и
\texttt{develop}, префиксы веток и префикс тегов v. Проверил, что все
ветки появились и запушил всё на GitHub. (см.
\hyperlink{fig-020}{рис.20}).

\begin{figure}

\centering{

\includegraphics[width=0.7\linewidth,height=\textheight,keepaspectratio]{image/20.png}

}

\caption{\label{fig-020}Выполнение git flow init, настройка master и
develop, префиксов веток и префиксв тегов v. Проверка, появления веток и
отправка на GitHub}

\end{figure}%
\end{frame}

\section{Выводы}\label{ux432ux44bux432ux43eux434ux44b}

\begin{frame}{Выводы}
В ходе данной лабораторной работы были получены навыки правильной работы
с репозиториями git.
\end{frame}




\end{document}
